\section{Future Work}
\label{sec:future}

Several concrete extensions would improve either the design or the
credibility of the evaluation. We list them in the order we consider
useful; the first is no longer future work, and we keep the section to
record what it was and what came of it.

\subsection{A Flat Temporal Index for Unbounded Ranges --- Implemented}
\label{sec:future:global-index}

An earlier revision of this section proposed a flat, sorted temporal
index alongside the year-bucketed one, so that an open-ended range query
could binary-search to its starting point instead of degrading to a
full-bucket fallback. That proposal has since been implemented
(\texttt{sfdb.sheaf.indexes.TemporalIndex.facts\_ending\_after} and its
symmetric counterpart \texttt{facts\_starting\_before}) and benchmarked;
Section~\ref{sec:discussion:temporal-unbounded} tells the full story. We
keep the entry, marked implemented, because a future-work item that
turned out to be an afternoon's change rather than a research programme
is itself evidence about how easily this class of gap is mistaken for a
design ceiling.

\subsection{Real-World Datasets and Standard Benchmarks --- Done}
\label{sec:future:real-data}

An earlier revision of this section named the evaluation's reliance on
synthetic data as a gap. Section~\ref{sec:eval:wikidata} has since run
the full benchmark against a genuine, organic dataset (Wikidata
office-holder facts we did not construct), and
Section~\ref{sec:eval:lubm} against LUBM~\cite{guo2005lubm}, the
standard benchmark workload for RDF stores, scaled up to
\lubmlargescale{} facts. Both moved the evaluation from
speculation to evidence, in different ways. LOOKUP, GLOBAL, and CONTEXT
hold on both; TEMPORAL, the one query class this paper's synthetic
benchmark already reported as sensitive to the storage-layer control, turns
out to be genuinely sensitive to the dataset itself too --- it beats
both in-house baselines on Wikidata and sits at roughly parity with
them on LUBM's two largest scales, neither matching the synthetic
result nor each other. LUBM's topology (many large, mutually unrelated
top-level contexts, unlike either the synthetic generator's balanced
tree or Wikidata's one-populous-country-among-many shape) also surfaced
a genuine performance bug in SFDB's own CONTEXT resolution
(Section~\ref{sec:discussion:threats}'s tenth finding), fixed and
re-verified across the entire evaluation. What remains is breadth on
the dataset side (DBpedia, a different corner of Wikidata) and, on the
standard-benchmark side, running LUBM against an external store rather
than only this project's own three in-house engines, and at the much
larger scales published LUBM results use --- see
Section~\ref{sec:future:jena}.

\subsection{Integration With Production RDF Stores --- Done}
\label{sec:future:jena}

An earlier revision of this section named a comparison against a real,
disk-backed production RDF store as ``the single highest-leverage item
on this list,'' since every simulated reviewer persona this project
consulted (\texttt{research/reviewer\_simulation.md}) raised it first.
That comparison has since been run three times, against three
independently
engineered stores: Section~\ref{sec:eval:jena} reports
all five query classes against a real Apache Jena TDB2 store (accessed
over HTTP via a Fuseki server, exactly as an external client would),
Section~\ref{sec:eval:virtuoso} reports the identical methodology
against a real OpenLink Virtuoso store, and Section~\ref{sec:eval:graphdb}
reports it again against a real Ontotext GraphDB store, at
every scale this paper measures, with identical result counts on every
row against all three. The finding is a genuine crossover rather than a clean
win, and it reproduces independently across all three stores: SFDB
remains faster than all three on LOOKUP, GLOBAL, and CONTEXT at every scale, while all
three overtake SFDB on TEMPORAL from $10^4$ facts onward. We
regard this as a more credible and more informative result than a
one-sided win would have been, and the agreement between three unrelated
optimisers, from three different organisations, as strong evidence the pattern is a
property of mature production query engines generally rather than one vendor's
implementation --- which replaces the earlier framing of
this item as unstarted future work, the question the second revision of
this section posed of whether a two-store agreement would hold for a
third, and the original question of whether a single comparison would
generalise at all.

What remains open: three stores is not an exhaustive survey of
production RDF engines, so a fourth or fifth comparison (e.g. Blazegraph
or Stardog, or a labeled property graph store such as Neo4j, which
would additionally require a schema translation rather than only a new
adapter --- Section~\ref{sec:related-work:lpg} discusses why) would
strengthen the generalisation claim further, though
with diminishing marginal value at this point. Standard benchmark
workloads have since also been run against all three external stores,
not only against this project's own in-house engines: LUBM~\cite{guo2005lubm}
(Section~\ref{sec:eval:lubm-external}) reproduces the identical
LOOKUP/GLOBAL/CONTEXT-robust, TEMPORAL-crossover pattern against Jena,
Virtuoso, and GraphDB, at scales up to \lubmjenalargescale{} facts, with
both the dataset and the comparison engine varied at once --- the one
axis-combination no other comparison in this paper tests. What remains
is scale (LUBM has not been run at the much larger scales published LUBM
results use) and breadth: BSBM~\cite{bizer2009bsbm} and
WatDiv~\cite{aluc2014watdiv} have not been tried at all.
That remains future work, but with the core "have we ever measured
against a real store" question answered, three times over, and "have we
ever measured against a standard workload, against a real store, varying
both at once" question also answered, three times over.

\subsection{An Incidence-Algebra Query Layer}
\label{sec:future:incidence}

Our survey of neighbouring formalisms identified incidence algebras over
locally finite posets~\cite{rota1964moebius} as the closest mathematical
neighbour to the context-poset design
(Section~\ref{sec:related-work}): the context poset is exactly the kind
of structure on which M{\"o}bius-inversion-style aggregation is defined,
which suggests a principled query layer for roll-up and drill-down
across context hierarchies. We deliberately did not pursue this here ---
it is a follow-up contribution, not a missing piece of the present
evaluation --- but it is the extension we consider most likely to add
new query capability rather than only new speed.

\subsection{Incremental Restriction and Global-Section Maintenance}
\label{sec:future:incremental}

Restriction maps are computed on demand and cached. Incremental
maintenance---updating cached restriction maps and the global-section
cache on insert rather than invalidating them---would improve query
latency in workloads that mix inserts and queries.

\subsection{Concurrent and Distributed Execution}
\label{sec:future:distributed}

Global-section computation is embarrassingly parallel across independent
context sub-trees. A parallel implementation would remove the single-node
scalability ceiling. A distributed implementation would additionally
require a partitioning strategy for the context poset and an inter-node
restriction-map protocol.

\subsection{Compression of Contextually Redundant Data}
\label{sec:future:compression}

Section~\ref{sec:eval:memory} reports that SFDB uses \rangememkg{}
more resident memory per fact than the KG baseline, because each fact is
registered in roughly eight to ten open sets plus four auxiliary indexes.
Sections that share a common context or that agree under restriction on a
large sub-tree admit context-aware compression schemes that are not
available to a flat triple store: rather than storing a full section
reference per open set, a fact could store one canonical reference plus a
small bitset of the open sets it belongs to, reconstructed on read.
Realising this would directly reduce the memory overhead identified in
Section~\ref{sec:eval:memory}, which is currently the paper's primary
reported cost.

\subsection{Hybrid Storage --- Revisited}
\label{sec:future:hybrid}

An earlier version of this section proposed a hybrid system that would
store facts under both representations and let the optimiser choose the
access path per query, on the premise that this would combine SFDB's
LOOKUP advantage with a GLOBAL advantage that belonged to the KG baseline.
That specific premise no longer holds: Section~\ref{sec:evaluation} reports SFDB
faster than the KG baseline on GLOBAL as well as LOOKUP. The storage-layer
control did surface a case where a second, reified access path might
help: on
bounded TEMPORAL, the SFDB/KG-mem ratio straddles parity across scale
rather than settling on either side of it (\rangetemporalkgmem{},
Section~\ref{sec:eval:control}) --- a control that itself does not
persist to disk sometimes beats SFDB and sometimes does not, with no
consistent trend across scale that a hybrid design could reliably
target. We do not prioritise building it here for exactly that reason:
the case for a second representation would rest on beating a control
by a margin large and consistent enough to be worth the engineering,
and what we have instead is a ratio that crosses parity in both
directions depending on scale, at absolute latencies in the low
single-digit milliseconds where measurement noise is itself a live
concern (Section~\ref{sec:eval:insert} treats the same regime cautiously
for insert timing).
SFDB also remains faster than the
\emph{raw} KG baseline on every TEMPORAL row measured
(Section~\ref{sec:eval:query}), so the case for a second representation
would rest entirely on the unstable KG-mem margin, not on anything
measured against a persisting store. Doubling the memory cost
identified in Section~\ref{sec:eval:memory} --- the paper's central
tradeoff --- to close a gap of that size is not an obviously good trade. A
hybrid design remains worth revisiting if a real workload shows the gap
widening at scale, which this evaluation's four data points do not yet
show one way or the other.
